\section{Background and Motivation}

\begin{frame}
	\frametitle{Double Pendulum}
	
	\begin{figure}
		\includegraphics[width=0.3\linewidth]{double_pendulum.png}
		\caption{Double Pendulum.}
	\end{figure}

    \vspace{1em}

    For the sake of simplicity, we refer to the state as $q = [\theta_1, \theta_2]^T$.


\end{frame}

\begin{frame}
    \frametitle{Double Pendulum Equations}
    \footnotesize % Reduce the font size for better readability

    %The state is defined as $q = [\theta_1, \theta_2]^T$.
    
    \vspace{1em}
    
    \textbf{Inertia Matrix} $M(q)$
    \begin{equation}
        \begin{bmatrix}
            (m_1{+}m_2)l_1^2 + m_2l_2^2 + 2m_2l_1l_2 \cos(q_2) & m_2l_2^2 + m_2l_1l_2 \cos(q_2) \\
            m_2l_2^2 + m_2l_1l_2 \cos(q_2) & m_2l_2^2
        \end{bmatrix}
    \end{equation}

    \vspace{1em}
    
    \textbf{Gravity Vector} $G(q)$
    \begin{equation}
        G(q) = \begin{bmatrix}
            (m_1{+}m_2)g l_1 \sin(q_1) + m_2 g l_2 \sin(q_1{+}q_2) \\
            m_2 g l_2 \sin(q_1{+}q_2)
        \end{bmatrix}
    \end{equation}
    
    % 3. Damping & Coriolis
    \vspace{1em}
    \begin{columns}[T]
        \begin{column}{0.4\textwidth}
            \textbf{Damping} $D$
            \begin{equation}
                D =
                \begin{bmatrix}
                    b_1 & 0 \\
                    0 & b_2
                \end{bmatrix}
            \end{equation}
        \end{column}
        
        \begin{column}{0.57\textwidth}
            \textbf{Coriolis Vector} 
            \begin{equation}
                \mathbf{c}(\mathbf{q}, \dot{\mathbf{q}}) = 
                \begin{bmatrix}
                    -m_2 l_1 l_2 s_2 (2\dot{\theta}_1 + \dot{\theta}_2)\dot{\theta}_2 \\
                    m_2 l_1 l_2 s_2 \dot{\theta}_1^2
                \end{bmatrix}
            \end{equation}
        \end{column}
    \end{columns}

    \vspace{1em}
    
    \textbf{Dynamics Equation}
    \begin{equation}
        M(q)\ddot{q} + \mathbf{c}(q, \dot{q}) + D\dot{q} + G(q) = u
    \end{equation}

   
\end{frame}

\begin{frame}
    \frametitle{Motivation}

    \begin{itemize}
        \item Real physical systems are chaotic and non-linear
        \item Parameter estimation with sensor is expensive and noisy
        \item ...
        \item ...
    \end{itemize}

\end{frame}


